% Nouveau EGA V (cote 25): archive cote 25, site page 167 (PDF page 168).
% English translation of the checked French annotated master.
% AI-translated by GPT-6 Astra (OpenAI, ChatGPT Pro mode).
\scanpage{25/167}
\mspage{\hand{2}}

\[ v^1 : E\hand{\otimes}\underline{O}_{V^1} \hand{\longrightarrow} \hand{\underline{\Omega}}^1_{X/Y}\hand{\otimes}\underline{O}_{V^1} \]
\var{\xout{in terms of which}}{}{typewriter strike-out}{ega25_25_167-v0001} (itself deduced from the homomorphism of modules
locally free on $V^o$
\[ v : \underline{E}\hand{\otimes}\underline{O}_{V^o} \hand{\longrightarrow} \hand{\underline{\Omega}}^1_{X/Y}\hand{\otimes} \underline{O}_{V^o} \quad , \]
in terms of which $V^1$ was defined as \var{\inserted[typed]{the}}{the}{insertion}{ega25_25_167-v0002} subprescheme of $V^o$ on which
the rank of v is strictly \hand{$\leqslant$} m-1 = rank $E\hand{\otimes}\underline{O}_{V^o}$ - 1). For the definition
of the first isomorphism, cf page 5; for that of the second, df the differential
supplements (attached) to ``variation of the rank of a homomorphism of
loc free Modules''. We therefore obtain an exact sequence:
\[ 0 \hand{\to} \underline{E}\hand{\otimes}\underline{O}_{V^1} \hand{\xrightarrow{\ \textcolor{black}{i}\ }} \underline{N}_{V^1/X} \hand{\xrightarrow{\ \ \textcolor{black}{j}\ \ }} \underline{\mathrm{Hom}}_{\underline{O}_{V^1}}(W,K) \hand{\longrightarrow} 0 \qquad . \]
One must then verify \var{\deleted{\var{\xout{1\textsuperscript{o}) that}}{}{typewriter strike-out}{ega25_25_167-v0004}}}{}{deletion}{ega25_25_167-v0003}
1\textsuperscript{o}) that
\[ \hand{\mu}i : \underline{E}\hand{\otimes}\underline{O}_{V^1} \hand{\longrightarrow} \hand{\underline{\Omega}}^1_{X/Y}\hand{\otimes}\underline{O}_{V^1} \]
is precisely $v^1$, which then implies, by the hypothesis that
$v^1$ is strictly of rank m-1, that \var{\xout{V(u)}}{}{typewriter strike-out}{ega25_25_167-v0005} $\hand{\mu}i$ defines, by passage to the
quotient,
\[ \hand{/\mu'} : \quad \underline{\mathrm{Hom}}\,(W,K) \hand{\longrightarrow} W = \mathrm{Coker}\ v^1 \ , \]
and that V(u) is also the subprescheme of $V^1$ which makes $\hand{\mu'}$ strictly
of rank \hand{\unsure{$<|\leqslant$}}rank W \var{\deleted{\unsure{$-$}1}}{}{deletion}{ega25_25_167-v0006}; on the other hand, one \var{\xout{can also interpret}}{}{typewriter strike-out}{ega25_25_167-v0007} must
verify

2\textsuperscript{o}) that $\hand{\mu'}$ is precisely the homomorphism deduced from the quadratic
form \var{\xout{on}}{}{typewriter strike-out}{ega25_25_167-v0008} M on W with values in K, at the bottom of page 15. It follows
at once that V($\hand{\mu'}$) = V(\unsure{D}) \var{\xout{x}}{}{typewriter strike-out}{ega25_25_167-v0009}, where D is the discriminant of M, which
completes the proof.

The verification of 1\textsuperscript{o}) and 2\textsuperscript{o}) (where perhaps a sign
must be introduced?) can only be essentially mechanical, and is left to
Dieudonné.

\mspage{\hand{167}}
\tnote{25-167-1}{The source's df before the differential supplements is retained literally; loc means locally. Strictly followed by a less-than-or-equal sign, the alternation of V(u) with V(a) on the preceding page, the slash before the displayed mu-prime, and the unverified possible sign remain as written. The canceled French fragment qui\unsure{t|l}e has no independent English lexical content. The source's méchanique and diffétentiels do not produce English misspellings.}
\tnote{25-167-2}{The canceled peut aussi s’interpréter is rendered can also interpret after English one; the live continuation is one must verify. The French apostrophe after il has no English counterpart.}
